\scanpage{90}
\diagram{\begin{tikzcd}[ampersand replacement=\&,row sep=1pt,column sep=2pt]
\& \parbox[t]{.85\linewidth}{a) pro\deleted{\unsure{tore}}\inserted{groupe de t.m.} $T$ sur $Q$, avec $f\in T(Q)$ engendrant $T$ (Zar)}\\
\Updownarrow \& {}\\
\& \parbox[t]{.85\linewidth}{b) Sous-groupe \deleted{de} $M$ de $\bar Q^*$, stable par $\operatorname{Gal}(\bar Q/Q)$}\\
\unsure{\Vert|\Updownarrow} \& {}\\
\& \parbox[t]{.85\linewidth}{c) Sous-catégorie str. pleine $\mathcal M'$ de la catégorie des vectoriels sur $Q$ avec
automorphisme \inserted{semi-simple} [i.e. de la catégorie des $Q[X]$ modules
\inserted{semi-simples} de dim finie sur $Q$] stable par $\otimes$, $\oplus$, sous objet}
\end{tikzcd}}

\deleted{\unsure{$a)\Rightarrow b)$}} Soit donné un groupe de t.m. $T$, avec ``générateur''
$f\in T(Q)$, et soit donné $G_m\xrightarrow{j}T$, ce qui équivaut
à se donner $j^*:M\longrightarrow\mathbb Z$, \inserted{invariant par $\operatorname{Gal}(\bar Q/Q)$}, ou une \unsure{graduation}\ \inserted{fonctorielle} sur
les objets de $\mathcal M$. \deleted{\unsure{Si de plus, il existe un}}
\inserted{\unsure{i.e.} \deleted{\unsure{$\varepsilon^*$}\illegible{2}}}
i.e. soit donné un objet $\Lambda$ de $\mathcal M$ \inserted{de rang 1} de \unsure{poids|puissance|disons} \unsure{,}
\inserted{de $\varepsilon:T\longrightarrow G_m$, \deleted{tel que} $\varepsilon\circ j(\lambda)=\lambda^{\alpha}$,}
i.e. $j^*(\varepsilon^*)=\alpha$,
\inserted{i.e. $\Lambda$ de degré $\alpha$.}
\deleted{\unsure{Posons|Pour} $\unsure{\varepsilon^*|\varepsilon} \deleted{\unsure{(f)}}=q$}
\deleted{\unsure{et} \unsure{$q^*$} de $|\varepsilon(f)|$}
\margin{$q=\varepsilon^*=\varepsilon(f)\in M\cap Q^*$}
% circled margin note, at left of the overwriting

\diagram{\begin{tikzcd}[ampersand replacement=\&,row sep=1pt,column sep=2pt]
\& \parbox[t]{.84\linewidth}{a) $T^\circ/\operatorname{Im}j$ est compact.}\\
\Downarrow \& {}\\
\& \parbox[t]{.84\linewidth}{b) \inserted{$q>0$ et} Pour $\forall$ $m\in M$, on a $|m|=\deleted{\unsure{j(f)}}q^{j^*(m)/\alpha}$.}\\
\Updownarrow \& {}\\
\& \parbox[t]{.84\linewidth}{c) \inserted{$q>0$ et} Pour tout $(E,f)\in\mathcal M$, si $E$ est de degré $i$,
les valeurs propres de $f$ dans $E_{\mathrm C}$ sont de val. absolue $q^{i/\alpha}$.}
\end{tikzcd}}

En effet, les valeurs propres des \deleted{$(E,f)\in$\unsure{$f$|$\mathcal M$},} $f$, pour $(E,f)\in\mathcal M$
de degré $i$, sont les $m\in M$ avec $\unsure{j(m)|j^*(m)}=i$. Donc
\deleted{\unsure{b signifie que}} c) et b) sont équivalentes.

\deleted{[Ils impliquent a), car $T'=T^\circ/\operatorname{Im}j$ est un tore \deleted{\unsure{ayant}}
- \unsure{possédant} $f'$ ayant la vertu que pour \unsure{tte}
représentation dans un $E$, les valeurs propres de
$f'$ opérant sur $E_{\mathrm C}$ sont de val. abs. $1$. \deleted{\unsure{Or il nous}}
\inserted{Or pour un} tore opérant, \inserted{\unsure{de simple|de simple-}} un ss-groupe fermé de $\operatorname{GL}(E_{\mathrm C})$ \unsure{qui est}, \deleted{\unsure{un tel}} est
\deleted{\unsure{par conséquent, un tore de} $\operatorname{GL}$ \unsure{dont les} $^\circ$ \unsure{de} $\mathbb R$ \unsure{opérant}}
\deleted{\unsure{sur} $E_{\mathrm C}$ sont de val. absolue $1$.}
\inserted{cpt.} Comme $T'_{\mathbb R}$ est un ss-groupe fermé de ce
dernier, cqfd \unsure{précisément}.]}
% archive sheet number: 89

